\documentclass[12pt]{article}
\usepackage{amsmath, amssymb, geometry}
\geometry{a4paper, margin=1in}

\title{Practice Exercise 1 Solutions}
\author{}
\date{}

\begin{document}

\maketitle

\section*{Exercise 1: Labor Productivity Growth}
\textbf{Solution:}
\begin{enumerate}
    \item \textbf{Labor Productivity Before Automation:}
    \[ \text{Labor Productivity} = \frac{\text{Output}}{\text{Number of Workers}} = \frac{10,000}{100} = 100 \text{ units/worker/year} \]
    
    \textbf{Labor Productivity After Automation:}
    \[ \text{Labor Productivity} = \frac{25,000}{50} = 500 \text{ units/worker/year} \]
    
    \item \textbf{Percentage Increase in Labor Productivity:}
    \[ \text{Percentage Increase} = \frac{500 - 100}{100} \times 100 = 400\% \]
    
    \item \textbf{Discussion:}
    The 400\% increase in labor productivity suggests significant efficiency gains. However, the reduction in workers may lead to unemployment unless these workers are reskilled or absorbed into other roles. Wages might increase for remaining workers due to higher productivity.
\end{enumerate}

\section*{Exercise 2: The Impact of Automation on Employment}
\textbf{Solution:}
\begin{enumerate}
    \item \textbf{Workers Employed Before Automation:}
    \[ L = 100 - 2w \]
    Substituting $w = 20$:
    \[ L = 100 - 2(20) = 60 \text{ workers} \]
    
    \textbf{Workers Employed After Automation:}
    \[ L = 80 - 2w \]
    Substituting $w = 20$:
    \[ L = 80 - 2(20) = 40 \text{ workers} \]
    
    \item \textbf{Wage at Which No Workers Are Employed:}
    \[ L = 0 \]
    \textbf{Before Automation:}
    \[ 0 = 100 - 2w \implies w = 50 \text{ (\$50/hour)} \]
    \textbf{After Automation:}
    \[ 0 = 80 - 2w \implies w = 40 \text{ (\$40/hour)} \]
    
    \item \textbf{Discussion:}
    Automation reduces the number of workers employed at the same wage and lowers the maximum wage at which employment exists. This highlights potential job displacement caused by automation.
\end{enumerate}

\section*{Exercise 3: Skill Premium and Automation}
\textbf{Solution:}
\begin{enumerate}
    \item \textbf{Percentage Change in Wages:}
    \begin{align*}
    \text{High-Skilled Workers:} & \quad \frac{40 - 30}{30} \times 100 = 33.33\% \\
    \text{Low-Skilled Workers:} & \quad \frac{15 - 20}{20} \times 100 = -25\%
    \end{align*}
    
    \item \textbf{Total Wage Bill:}
    \textbf{Before Automation:}
    \[ \text{Wage Bill} = (30 \times 50) + (20 \times 100) = 1,500 + 2,000 = \$3,500 \]
    
    \textbf{After Automation:}
    \[ \text{Wage Bill} = (40 \times 50) + (15 \times 100) = 2,000 + 1,500 = \$3,500 \]
    
    \item \textbf{Discussion:}
    Despite unchanged total wage bill, income inequality rises as high-skilled workers benefit while low-skilled workers face wage reductions. This reflects the skill-biased nature of technological advancements.
\end{enumerate}

\section*{Exercise 4: Technological Adoption and Output Growth}
\textbf{Solution:}
\begin{enumerate}
    \item \textbf{Percentage Increase in Output:}
    \[ \text{Percentage Increase} = \frac{80,000 - 50,000}{50,000} \times 100 = 60\% \]
    
    \item \textbf{Change in Total Revenue:}
    \[ \text{Initial Revenue} = 50,000 \times 25 = \$1,250,000 \]
    \[ \text{New Revenue} = 80,000 \times 25 = \$2,000,000 \]
    \[ \text{Change in Revenue} = 2,000,000 - 1,250,000 = \$750,000 \]
    
    \item \textbf{Discussion:}
    The increase in output and revenue demonstrates how automation can enhance competitiveness and market share. However, cost savings might pressure competitors to adopt similar technologies, driving further automation.
\end{enumerate}


\section*{Exercise 5: Job Polarization}
\textbf{Scenario:} In an economy, job distribution shifts due to automation. The initial job shares across sectors are as follows:
\begin{itemize}
    \item Low-skill jobs: 40\%
    \item Middle-skill jobs: 40\%
    \item High-skill jobs: 20\%
\end{itemize}

After automation, the shares change:
\begin{itemize}
    \item Low-skill jobs: 45\%
    \item Middle-skill jobs: 25\%
    \item High-skill jobs: 30\%
\end{itemize}

\textbf{Questions:}
\begin{enumerate}
    \item Calculate the percentage point change in job shares for each skill level.
    \item If the total labor force is 10 million, how many jobs are there in each category before and after automation?
    \item Discuss how automation leads to job polarization and its implications for income inequality.
\end{enumerate}

\textbf{Solutions:}
\begin{enumerate}
    \item Percentage point change in job shares:
    \begin{itemize}
        \item Low-skill jobs: $45\% - 40\% = 5$ percentage points (increase).
        \item Middle-skill jobs: $25\% - 40\% = -15$ percentage points (decrease).
        \item High-skill jobs: $30\% - 20\% = 10$ percentage points (increase).
    \end{itemize}
    \item Number of jobs in each category:
    \begin{itemize}
        \item Before automation:
        \begin{itemize}
            \item Low-skill jobs: $40\% \times 10,000,000 = 4,000,000$.
            \item Middle-skill jobs: $40\% \times 10,000,000 = 4,000,000$.
            \item High-skill jobs: $20\% \times 10,000,000 = 2,000,000$.
        \end{itemize}
        \item After automation:
        \begin{itemize}
            \item Low-skill jobs: $45\% \times 10,000,000 = 4,500,000$.
            \item Middle-skill jobs: $25\% \times 10,000,000 = 2,500,000$.
            \item High-skill jobs: $30\% \times 10,000,000 = 3,000,000$.
        \end{itemize}
    \end{itemize}
    \item Automation often leads to job polarization by increasing demand for both high-skill and low-skill jobs while reducing demand for middle-skill jobs. This can widen income inequality as high-skill jobs command higher wages, while low-skill workers face stagnant or declining wages.
\end{enumerate}

\section*{Exercise 6: Automation's Effect on GDP Growth}
\textbf{Scenario:} A country’s GDP is modeled as:
\[ GDP = A \times L^{0.6} \times K^{0.4} \]
where $A$ is the technology level, $L$ is labor input, and $K$ is capital input.

Initially:
\begin{itemize}
    \item $A = 1$, $L = 1,000$, $K = 500$.
\end{itemize}

After automation:
\begin{itemize}
    \item $A = 1.2$, $L = 800$, $K = 600$.
\end{itemize}

\textbf{Solutions:}
\begin{enumerate}
    \item Calculate GDP:
    \begin{itemize}
        \item Before automation:
        \[ GDP = 1 \times 1,000^{0.6} \times 500^{0.4} \]
        \[ GDP = 1 \times 100^{0.6} \times 50^{0.4} = 1 \times 63.1 \times 18.8 \approx 1,187.28 \]
        \item After automation:
        \[ GDP = 1.2 \times 800^{0.6} \times 600^{0.4} \]
        \[ GDP = 1.2 \times 79.4 \times 22.1 \approx 2,111.01 \]
    \end{itemize}
    \item Percentage change in GDP:
    \[ \text{Percentage change} = \frac{2,111.01 - 1,187.28}{1,187.28} \times 100 \approx 77.8\% \]
    \item Technological advancements significantly increase GDP by improving productivity (higher $A$). While labor input decreases, higher capital and technological efficiency compensate, leading to economic growth. However, the reduction in labor input may create challenges such as unemployment or skill mismatches.
\end{enumerate}


\end{document}
